\documentclass[aps,pre,10pt,twocolumn,superscriptaddress,nofootinbib]{revtex4-2}
\usepackage{graphicx}
\usepackage{amsmath}
\usepackage{booktabs}
\usepackage[hidelinks]{hyperref}

% Working note, not part of the paper (13 October 2026). Published relatives to read before the next revision:
% order melted by injected energy, and its recovery, in electron-lattice models [BBM26], [BFC26], [YC24]
% (REFERENCES.bib; abstracts and parts only). [BBM26] separates melting that loses agreement between regions while
% local order survives from melting that destroys the local order; our local counts could say which kind each
% "melt" here is, and that reading has not been made. Vocabulary and sources:
% docs/reading/notes/2026-10-13_chern_group_vocabulary.md, sections 2, 5 and 6.

\begin{document}

\title{Concentrated energy does not fold space: seven pre-registered tests in a graph toy model of emergent
geometry}

\author{Emily Smith}
\email{emilysmithcreate@gmail.com}
\affiliation{Independent researcher, Charlottesville, Virginia, USA}

\date{DRAFT, 14 October 2026}

\begin{abstract}
The author's hypothesis holds that a black hole is a region of space re-curled toward the arrangement it came from,
paid for by concentrated energy. Combinatorial quantum gravity (CQG) reads it as a bubble of the random phase. We ask
a graph toy: CQG's energy with the coefficient $\lambda$ of its local term raised above 1 (only $\lambda = 1$ is CQG).
Flat space is given energy, and vertices are classified by a count of squares as folded or melted. Seven
pre-registered tests cover two, three and four directions; energy spread, packed or drained; points named or
interchangeable; and, in the last, a tie between directions that the author chose. Every overall verdict is a form of
MELTS or HEALS. One partial reading is an exception by the letter: at 36 vertices, where no complete fold exists, one
protocol reads as folding on three damaged replicas. Fold predictions recorded for the author failed in all seven;
ours failed in part in two and missed details in two more. In three directions at moderate energies the melted
states, read from the wiring, are a few flickering scars that heal when cooled. Under the author's tie, with twice
the price of curling a whole direction packed into one place, half the space ends broken and none of it curled.
Exact: no single move out of a flat sheet adds a square.
\end{abstract}

\maketitle

\section{What is asked}\label{sec:asked}

This is paper 12 of a series of short papers \cite{repo}. Companion papers study a curled torus that converts to a
flat one \cite{paper1}, what it leaves behind \cite{paper2}, and how curled directions open with six and eight links
\cite{paper5}. This paper asks the reverse. Can energy put into flat space curl it back up?

\emph{The author's claim.} A black hole is a closed region of space re-curled toward the curled arrangement that space
came from. The energy concentrated in it pays the curling cost. In open space energy spreads and space heals; in a
closed region space folds. Space itself may be stable: the return is paid for, not fallen into \cite{repo}.

\emph{The published reading.} In combinatorial quantum gravity (CQG) a black hole is a large bubble that has reverted
to the random phase \cite{T25}. That work reads the random phase as matter. The author does not.

The two readings differ in what concentrated energy does to flat space: curl it, or disorder it. A sealed simulation
can count both.

\emph{A benchmark.} In the author's wider picture the fold is the return leg of a cycle, and new space is born where
black holes collide. A published cousin is cosmological natural selection, in which black-hole singularities bounce
into new expanding universes \cite{Smolin06}. It is a benchmark for what such a cycle must do; nothing here tests it
or draws support from it.

\emph{Conventions.} Each statement is one of four kinds, and we label it: \emph{exact} (arithmetic or complete
enumeration, at the sizes stated); \emph{measured} (a run registered before it ran, with its verdict label in
capitals); \emph{interpretation} (ours, unverified); or \emph{the author's} (the hypothesis's claim). Runs that were
not registered are marked exploratory. Only $\lambda = 1$ is CQG. Every test here is at $\lambda > 1$, two of them
with a further term, a tie between directions. We call this family the toy; it is not CQG.

\section{The toy, the sealed box and the reading rule}\label{sec:toy}

\emph{State and energy.} The state is a $2D$-regular two-sided (bipartite) graph on $N$ vertices under the hard-core
rule of CQG \cite{T25}. We use $D = 2$, 3 and 4: four, six and eight links per vertex, which we call two, three and
four directions. With $S$ squares (4-cycles), $S_e$ of them on link $e$, and $X$ surplus squares on over-full links,
\begin{equation}
\begin{aligned}
  H &= 16\left[\tfrac{1}{2}D(D-1)N - S\right] + 4\lambda X ,\\
  X &= \textstyle\sum_e (S_e - 2D + 2)_+ .
\end{aligned}
\label{eq:H}
\end{equation}
The flat $D$-torus has $H = 0$; we call it flat space. A direction of length four is called curled. \emph{Exact:} a
curled direction costs $a = 4(\lambda - 1)$ per vertex ($\epsilon$ in Ref.~\cite{paper1}), and the costs of several
curled directions add \cite{paper5,repo}. At four links Eq.~(\ref{eq:H}) is the published energy with $\lambda$ raised. The
six- and eight-link forms are derived by us from the published definitions.

\emph{Reproduction gate, stated once.} Our six-link code has not reproduced the published six-link curve (Fig.~3
of Ref.~\cite{T22}). A factor of two in the coupling is accounted for by the published equation; the remaining
difference is unsettled. No published eight-link curve exists to compare with. Run at four links, the same code
makes the four-link chain draw for draw, and its ladder of curled states is exact. Every six- and eight-link
statement below is about our code and carries this open gate \cite{repo}.

\emph{The reading rule.} The count $d(v)$ is the number of pairs of links at $v$ that close no square. It is $D$ in
flat space (the flat value), falls by one for each curled direction, and rises where squares are missing. A
vertex is \emph{folded} if $d < D$ and \emph{melted} if $d > D$; both count as damaged. A replica, read from its
final measurement block, is FOLDED if at least four vertices are damaged and folded ones are at least as many as
melted, MELTED if at least four are damaged and melted outnumber folded, and HEALED if fewer than four are damaged.
Four is the smallest curled object, one column. A cell (one choice of size, protocol and energy) takes the majority
outcome of its replicas. A test reads RE-CURLS if
some cell has a FOLDED majority, MELTS if none has and some cell has a MELTED majority, and HEALS if every cell is
HEALED. Two tests used other registered rules: T21 the mean folded share of the damage, and T47 shares of vertices
(Sec.~\ref{sec:four}). These labels are counts. The counts alone do not establish geometry, dimensionality
or a return to the random phase. Sections~\ref{sec:two} and \ref{sec:three} say what the MELTED states can be, and
what they were where their wiring was read.

\emph{The sealed box.} Stores of nonnegative energy (Creutz reservoirs) sit beside the graph \cite{Creutz83}. A
proposed switch (two links swap their endpoints on one side, keeping every degree) draws on one store, is refused
if the store cannot pay, and pays into it when it lowers the energy. $H$ plus the stores, plus whatever has been
drained, is conserved by construction. Conservation is bookkeeping; it does not establish
equilibrium. Three protocols differ in where the energy starts and which store a move may use. \emph{Spread}: a
shared bath of $2N$ stores with the whole budget in one, which any move may draw on. \emph{Packed}: one store per
vertex, the budget in the store of one vertex, and a move pays only from the store of the first vertex it involves;
or the budget wired into the graph as disorder near one vertex, with the bath empty. \emph{Drained}: the spread
protocol with every store scaled by $1 - \ell$ after each block of 100 sweeps. The cheapest move out of flat space
costs 32, 64 or 128 with four, six or eight links (Sec.~\ref{sec:exact}), so a single store holding less can never
pay for it.

\emph{Named or interchangeable points.} By default each vertex has its own name. In the interchangeable
(unlabeled) ensemble a graph gains weight in proportion to its number $A$ of side-preserving automorphisms, and
the ratio $A(G')/A(G)$ enters the acceptance of a move \cite{DQM25,repo}. This can favor symmetric arrangements. It
does not guarantee that the fully folded state outweighs every partly folded one, and it says nothing about how fast
a fold forms. It is a modeling choice, not derived here from quantum indistinguishability.

Table~\ref{tab:tests} lists the seven tests.

\begin{table*}[t]
\caption{The seven pre-registered tests. $D$ is the number of directions ($2D$ links per vertex). Verdict labels are
the registered ones. ``Inferred'' means inferred for the author by the assistant from her stated positions.}
\label{tab:tests}
\small
\setlength{\tabcolsep}{3pt}
\begin{tabular}{@{}lclcp{3.0cm}p{3.5cm}p{3.0cm}p{2.6cm}@{}}
\toprule
Test & $D$ & $N$ & $\lambda$ & Protocol; points & Verdict & Predicted for the author & Ours \\
\midrule
T21 & 2 & 36, 64 & 1.25 & spread; named and interchangeable & MELTS EITHER WAY (overall, $N = 64$); at
$N = 36$, bath, FOLDS WITH INTERCHANGEABLE POINTS by the letter & a closed region folds (inferred from her
statement): failed at 64 & MELTS EITHER WAY: held at 64, failed at 36 (bath) \\
T26 & 2 & 64, 144, 256 & 1.25 & packed; named, and both at 64 & MELTS & RE-CURLS (inferred, then confirmed by her):
failed & MELTS: held \\
T27 & 2 & 144 & 1.25 & spread, then drained; named & STAYS MELTED & FOLDS BEFORE IT FLATTENS (inferred): failed &
STAYS MELTED: held \\
T34 & 3 & 216, 512 & 1.25 & packed and spread; named & MELTS & RE-CURLS (her words): failed & MELTS: held \\
T42 & 3 & 512 & 1.02 & spread; interchangeable, with a named control & HEALS; control MELTS & RE-CURLS, control
MELTS (inferred): failed; held & MELTS in both: failed for interchangeable \\
T47 B & 4 & 2,304 & 1.30 & packed; named; a tie table (Sec.~\ref{sec:four}) & HEALS; control HEALS & FRONT, control HEALS
(inferred, then confirmed by her): failed; held & HEALS in both: held \\
T57 & 3 & 216, 512 & 1.10, 1.25 & packed and spread; named; the author's tie (Sec.~\ref{sec:tie}), and untied &
MELTS, tied and untied & RE-CURLS with the tie (her words): failed & MELTS or HEALS, tied and untied: held \\
\bottomrule
\end{tabular}
\end{table*}

\section{Two directions}\label{sec:two}

All runs in this section are on flat four-link tori at $\lambda = 1.25$ unless stated.

\emph{An earlier exploratory run.} A flat sheet of 144 or 160 vertices was given 0.5 to 16 units per vertex in a single
store, at $\lambda = 1.05$ to 2.0. Melted vertices outnumbered folded ones at every budget and every $\lambda$. The
prediction written before the run, that small budgets buy folds because a fold is the cheaper damage, failed. As
the author pointed out, heating a sheet at equilibrium randomizes it by construction; the run compares only two ways
of holding the same energy, and folding lost \cite{repo}.

\emph{Spread (T21), measured.} $N = 36$ and 64; budgets of 2, 4, 8 and 16 per vertex through the bath and 0.5 to 4
through a single store; six replicas per cell; named and interchangeable points; 4,000 sweeps. At $N = 64$, where
complete folds fit (four 4-cubes, or the $16 \times 4$ curled torus), the mean folded share of the damage under the
bath is at most 0.27 with interchangeable points and 0.18 with named ones, at budget 2, and falls to 0.02 at 16. The
registered overall reading is the one at $N = 64$ under the bath: MELTS EITHER WAY. The single store reads MELTS
EITHER WAY at both sizes. At $N = 36$, where no complete fold exists, the bath reads FOLDS WITH INTERCHANGEABLE
POINTS by the letter: 0.61 against a named 0.43 at budget 2, from three damaged replicas, with the order reversed at
budget 4 (0.25 against 0.40). The author's statement that a closed region folds corresponds to FOLDS WITH
INTERCHANGEABLE POINTS; she had made no separate prediction. It failed at $N = 64$. Ours, MELTS EITHER WAY at both sizes,
held at 64 and failed at 36 under the bath.

\emph{The author's objection to spreading:} a black hole is energy packed into one place.

\emph{Packed (T26), measured.} Twenty-seven cells: 32 to 192 units in one vertex's store at $N = 144$ and 256; 288
and 576 units over nine stores at $N = 144$; a disordered patch of 32 to 128 units in a cold shared bath at $N =
144$ and 256, and with per-vertex stores at $N = 144$; and at $N = 64$ a patch of 32 or 64 units with named and with
interchangeable points. Twenty replicas per cell and 50,000 sweeps (ten and 4,000 at $N = 64$). With the budget in
one store, 19 or 20 of 20 replicas are MELTED in every cell; over nine stores, 20 of 20. A patch in a shared cold
bath: 15 to 20 of 20 MELTED, the rest HEALED, at most one FOLDED. A patch with per-vertex stores: 20 of 20 MELTED.
At $N = 64$ the named patch is MELTED in 7 of 10 at both budgets, with 3 FOLDED at 32; the interchangeable one is
HEALED in 10 of 10 at 32 and MELTED in 7 of 10 at 64. No cell has a FOLDED majority. The largest folded piece in any
final state is 11 vertices, and at most 1 replica in 20 melts first and folds later. Verdict: MELTS, with 26 cells
MELTED, one HEALED and none FOLDED. The author's prediction, RE-CURLS with melt-then-fold common, was inferred by
the assistant and confirmed by her before any cell was read as a verdict. It failed. Ours, MELTS, held, with two
misses: we expected the patch in a shared cold bath to heal at 144 and 256 vertices, and it mostly did not; and we
expected the two kinds of point to agree at $N = 64$, and they differed at 32 units.

\emph{What MELTED can mean at the smallest budgets.} \emph{Exact:} the cheapest move out of a flat sheet loses two
squares and costs 32. It leaves eight vertices at $d = 3$, as the registration states and its disclosed trial run
showed. By the rule, a replica that buys that one move and nothing more has eight damaged vertices, none folded,
and reads MELTED. So at the smallest packed budgets the label can mean one broken-square defect, not a disordered
region. The two-direction end states were classified by the count; their wiring was not read, as it was in three
directions (Sec.~\ref{sec:three}).

\emph{Drained (T27), measured.} $N = 144$; the spread protocol at 2, 4 and 8 per vertex; drained at
$\ell = 0.001$, 0.01, 0.1 and 1.0 per block, the last a quench; a sealed control; twelve replicas per cell; 30,000
sweeps. All fifteen cells, the control included, have a MELTED majority. Of the 144 draining replicas 134 are MELTED
and 10 FOLDED. No replica melts first and folds later, and the largest folded piece anywhere is 16 vertices. The
quench freezes the disorder as it stands, and the slow drains do not anneal it within the run. Verdict: STAYS MELTED.
The prediction inferred for the author by the assistant, FOLDS BEFORE IT FLATTENS, failed; the record does not
show her confirming it. Ours, STAYS MELTED, held as a verdict; we had expected the slow drains to flatten the
sheet, and they did not.

\emph{Left alone (exploratory).} With a chain that also allows triangles and pentagons, run at fixed coupling $g$
and not sealed, a flat sheet at $g = 1$, 2 and 4 accepted no move at all, and at $g = 8$ it took up 7.9 units per
vertex as disorder. It did not fold. A third part of that run, an apparent jammed state, was withdrawn as an
artifact of a slow chain \cite{repo}.

\section{Three directions}\label{sec:three}

\emph{The author's mechanism.} The collapse begins on one direction. Once one direction has curled the rest follow,
so the grouped directions collapse together. It has two halves: the walls (each fold makes the next easier) and the
counting (with interchangeable points the fully curled, symmetric state weighs more). Both halves are in the toy's
arithmetic (Sec.~\ref{sec:exact}). T34 tests the first and T42 the second.

\emph{Walls, named points (T34), measured.} Flat $6 \times 6 \times 6$ ($N = 216$) and $8 \times 8 \times 8$ ($N =
512$) tori at $\lambda = 1.25$. Energies run from 64, the wall of flat space, to one and two
directions' worth of curling cost plus a wall (260 and 500 at $N = 216$; 560 and 1,060 at $N = 512$). The energy is
packed into one vertex's store or spread through $2N$ stores; eight replicas per cell; 50,000 sweeps. All 20 cells
have a MELTED majority. No replica ended with a single folded vertex, none reached a whole direction's worth of
folded vertices at any reading, and none melted first and folded later. Verdict: MELTS. The author's prediction,
RE-CURLS, one direction first and then the rest, failed. Ours, MELTS, held.

\emph{What those MELTED states are.} In an exploratory quench, with its expectation written first, all 80 final
graphs at $N = 216$ were run for 5,000 sweeps with only moves that do not raise the energy. All 80 came to rest on
exactly flat space, with no defect left. Read later from the wiring, with T42's below, the MELTED replicas
end with one to a few small scars, 12 to 46 vertices of 216 or 512, most of them one swapped pair of links,
flickering on and off. The largest excitation, about 150 of 512 vertices at the largest energy, healed back to such
scars. Nothing curled. The verdict stands under its rule but says less than its name: the energy excites flat
space, and the excitation heals \cite{repo}.

\emph{Counting, interchangeable points (T42), measured.} Flat $8 \times 8 \times 8$ at $\lambda = 1.02$, where
curling all three directions of the whole torus costs 123 units. The spread protocol with 64 to 2,048 units; eight
replicas per cell; 50,000 sweeps; and a control with the same chain and the automorphism factor off (named points).
Interchangeable points: HEALS, with 47 of 48 replicas HEALED and one MELTED at the largest energy. Named control:
MELTS, with a MELTED majority in five of six cells. The author's prediction, inferred for this design by the
assistant, was RE-CURLS with interchangeable points and MELTS in the control. The first half failed and the second
held. Ours, MELTS in both with no visible difference, failed for the interchangeable half.

\emph{The named control did not melt (a correction on the record).} Its MELTED replicas hold 12 to 42 vertices of
512 above the flat value at the last reading: one to a few scars, most of them 12 vertices holding 64 units,
one swapped pair of links, which is the cheapest move out of flat space. At the smaller energies the scar switches
on and off about 40 times a run and is present about three quarters of the time, so the last reading is a snapshot
of a flicker. At the largest energy up to about half the vertices are disordered early on, with either kind of
point, and heal back. At most 2 of 512 vertices ever fell below the flat value. Both treatments heal. With named
points a small scar keeps flickering; with interchangeable points it almost never forms. \emph{Interpretation (ours,
unverified):} a scar at each of its many possible places is a different named arrangement but one interchangeable
arrangement, so the scar has far fewer ways to exist. Here the counting works against the scar. It does not work for
the fold.

\section{Four directions}\label{sec:four}

\emph{The author's claim (26 September 2026).} Time curls behind the present: the present is a moving front, and
a state with one direction curled lies below flat four-direction space. The toy has no time. It can ask only whether
flat eight-link space re-curls from a local push when a one-curled state lies lower.

\emph{Measured (T47, part B).} Flat $6 \times 6 \times 8 \times 8$ ($N = 2{,}304$) at $\lambda = 1.30$, with a tie
table: an extra cost per vertex that depends on its number of open directions. Under the test table the one-curled
state lies 0.24 per vertex below flat space; under the control table flat space is lowest. The cheapest move out of
flat space costs 128 under both. One push of 128.01 units was placed in one vertex's store; three replicas per
table; 20,000 sweeps. The registered rule: MELTED if a quarter or more of the vertices are above $d = 4$, FRONT if a
quarter or more are at $d = 3$, HEALS if fewer than 2\% are anywhere but $d = 4$, STALLED otherwise. Verdict: HEALS
under both tables. At most 0.9\% of the vertices ever left $d = 4$, and the runs end with a damaged patch of about
20 vertices that holds the push and neither spreads nor curls. The author's prediction, inferred by the assistant
and confirmed by her before any result was read, was FRONT under the test table and HEALS under the control. The
first failed and the second held. Ours, HEALS under both, held. \emph{Interpretation (ours, unverified):} a patch of
20 vertices gains about $20 \times 0.24 \approx 5$ units from lying lower, far less than its boundary costs, so a re-curled region would have
to be large before it grew by itself. A larger or more concentrated push was not tested. The table's values came
from a fit to an order of openings that the author has since set aside \cite{repo}; that does not change what the
run tested.

\section{What is exact}\label{sec:exact}

\emph{The floor.} At four links there are $2N$ links and each square has four, so $\sum_e S_e = 4S$ and
$16(N - S) = 4\sum_e (2 - S_e)$. Eq.~(\ref{eq:H}) is then a sum over links,
\begin{equation}
  H = 4\sum_e \bigl[(2 - S_e) + \lambda\,(S_e - 2)_+\bigr] . \label{eq:Hlink}
\end{equation}
The hard-core rule limits a link $uv$ to three squares: each runs through a different other neighbour of $v$,
since two through the same $w$ would give $u$ and $w$ three common neighbours. The term is $4(2 - S_e) \geq 0$ for $S_e \leq 2$ and $4(\lambda - 1)$
for $S_e = 3$, so at $\lambda \geq 1$ every term is nonnegative and nothing lies below the flat sheet. Flat space
is stable there, not metastable. On the six- and eight-link ladders flat space is the lowest rung at $\lambda > 1$.

\emph{No first step.} Enumerating every switch the chain can make out of a perfect $10 \times 10$ or
$12 \times 12$ sheet: none adds a square; 14,900 and 34,128 lose squares; the cheapest costs 32. Out of a
disordered graph 5,127 switches add a square, and the cheapest of them lowers the energy by 75. So a fold cannot begin in
one step from a flat sheet. One order can still be built from another through a thin seam where squares are briefly
broken, as the opening of the curled torus shows in reverse \cite{paper1}. In the runs above the broken squares did
not go on to fold.

\emph{Price does not decide it.} At $\lambda = 1.25$ a square added to a flat sheet costs $16\lambda - 16 = 4$
(each of its four links gains a surplus square) and a square destroyed costs 16, so per square a fold is the cheaper damage by a factor of four. In the exploratory runs of
Sec.~\ref{sec:two} the folded share of the damage rose as $\lambda$ fell toward 1, to 23\% at most at $\lambda =
1.05$ against 10\% at most at 1.25, and never reached one half.

\emph{The wall of flat space.} The cheapest move out (32, 64 and 128 at four, six and eight links) loses 2, 4 and 8
squares and adds no surplus square, so it does not depend on $\lambda$.

\emph{The walls' half of the mechanism.} On the six-link ladder, on the tori and the gas of 6-cubes computed, the
cheapest single move out of a state with one, two and three curled directions costs $96 - 48\lambda$,
$96 - 64\lambda$ and $96 - 80\lambda$: 36, 16 and downhill at $\lambda = 1.25$. Two neighbouring rungs share the top
of the barrier between them, so the cost of a fold beyond its own curling cost is the wall out of the more-curled
rung, and each fold lowers the next one's excess by $16\lambda$. Nothing in the untied energy ties
the directions together. \emph{Interpretation (ours, unverified):} a single move need not be the whole barrier in
three directions, where the boundary of a curled patch can grow with the patch. That barrier is priced in a separate
registered calculation whose verdict is not yet read \cite{paper5,repo}.

\emph{The counting's half.} With interchangeable points the free energy of an arrangement is $H - g \ln A$. At
$N = 512$, $A = 12{,}288$ for flat $8 \times 8 \times 8$, 2,048 with one direction curled, 12,288 with two, and
$3.2 \times 10^{39}$ for the fully curled state, a gas of eight 6-cubes. The partial folds carry no more symmetry
than flat space; only the fully curled state gains. At $g = 1.5$ it beats flat space only for $\lambda < 1.020$
(1.013 at $g = 1.0$). Identical cubes can be exchanged, so the crossing rises slowly with their number: 1.031 at
$10^3$ cubes, 1.045 at $10^6$, and 1.25 only near $10^{52}$. This compares single arrangements and ignores their
thermal excitations. T42 ran at $\lambda = 1.02$, and nothing folded.

\section{The author's tie, and the test under it}\label{sec:tie}

\emph{The author's decision (9 October 2026).} The three directions are tied, ``all at the last.'' A vertex with $d$
of its three directions open pays an extra $f(d)$, with $f = (0, a, 2a, 0)$ for $d = 0, 1, 2, 3$. She regards the
untied runs as a stepping stone.

\emph{Exact under the tie.} Every rung of the six-link ladder sits $3a$ per vertex above flat space ($k \geq 1$
curled directions cost $ka$, and the tie adds $f(3 - k) = (3 - k)a$), so nothing is released until a vertex's last
direction opens. Read as a fold: curling one direction of flat space costs $3a$ per
vertex, the same as curling all three. The first fold pays the whole price and the later folds are free. The wall of
flat space stays 64 at every $\lambda$ computed, 1.05 to 1.25. As implemented, a vertex above the flat value pays
no tie, so the tie charges curling and not breaking.

\emph{The author's mechanism for the fold (her words, 9 October).} A high amount of energy concentrates in space,
then more arrives, until ``there is just too much and it has to go somewhere and the only mechanism it can go to is
curling macrodimensional space itself.''

\emph{Under the tie (T57), measured.} None of the six tests above used this tie. T57 is T34's protocol on flat
six-link tori of 216 and 512 vertices, with the tie and with an untied control, at $\lambda = 1.10$ and 1.25, packed
and spread, eight replicas per cell and 50,000 sweeps, at energies up to 3,200 units at $N = 512$: about twice what
curling one direction of the whole torus costs under the tie at $\lambda = 1.25$. The reading rule is T34's, with a
verdict for each combination of tie and $\lambda$. Verdict: MELTS in all four groups. Of 640 replicas 582 are MELTED, 58
HEALED and none FOLDED. No replica reached a whole direction's worth of folded vertices; the most vertices below the
flat value at any reading is 15 with the tie and 25 without. The author's prediction, as the registration reads
her words, was RE-CURLS with the tie at the largest energies. It failed. Ours, MELTS at $\lambda = 1.25$ and MELTS or
HEALS at 1.10, tied and untied alike, with no FOLDED replica, held.

\emph{Where the energy went.} Into broken squares. Up to about one direction's worth, a replica ends with a median of
12 damaged vertices, one swapped pair of links, as in T34. Packed into one store at one direction's worth and at
twice it, a median of 97 to 107 and then 255 to 263 of 512 vertices are damaged at the end, the same with the tie as
without. Spread through a bath the same energies leave 12 to 79. At the largest packed energies MELTED means what it
says by the count: half the torus above the flat value. Those end states' wiring has not been read.
\emph{Interpretation (ours, unverified):} the energy had another place to go. As implemented the tie charges a curled
vertex and not a broken one, so it makes breaking the cheaper of the two; and with or without it no single move out
of flat space adds a square. Two exact counts would separate these and are not done: the price of the smallest curled
patch in flat space against the same number of broken vertices, and the same under a table that charges damage. Such
a table would be a new rule, and it is the author's to choose \cite{repo}.

\emph{What a fold would need (interpretation, ours, unverified).} Before anything in the toy is called a black hole
it would have to show, in order: a stable, localized curled region inside open space; a reproducible threshold;
growth with stored energy; trapping or an analog of it; and a sensible large-scale limit. With four links one such
region is known, the relic left when a curled torus opens: one curled column, a local minimum costing $24\lambda -
16$ that lasts only in the cold \cite{paper2}. With six links none has been found: not among one- and two-switch
constructions (exact within that search, at $\lambda = 1.02$), and not as a remnant of heat damage in the quench of
Sec.~\ref{sec:three}. The toy also has no pull: at fixed wiring, defects that share no square do not interact
through the energy \cite{paper11}. Nothing in it gathers energy, so the energy is placed by hand. Where the entropy
attributed to a black hole would live in an ordered interior is unanswered.

\section{What this does not show}

It shows nothing about real black holes. It is a statement about a graph toy with no time and no pull, at 36 to 512
vertices (2,304 with eight links, on three replicas per table), over sweeps that are not a physical clock.

No registered test here is at $\lambda = 1$, so no verdict here is a result in CQG.

The packed protocol was not run with interchangeable points in three directions, and the tie runs with named points
only.

MELTED is a count. \emph{Interpretation (ours, unverified):} where the wiring was read, in three directions at
these energies, the toy produced neither the author's fold nor a lasting disordered bubble; flat space was excited
and healed. In two directions the largest spread budget of the exploratory run left every vertex melted by the
count, as equilibrium heating must. Whether that state is the random phase that
Ref.~\cite{T25} reads as a black hole's interior was not measured. A correction of the author's stands beside this:
a region returned to the random phase is not obviously the opposite of her picture, and may be a different account
of the interior.

Every six- and eight-link statement carries the open reproduction gate of Sec.~\ref{sec:toy}. No physicist has
reviewed this work.

\begin{acknowledgments}
Code, analysis and text were drafted with Claude (Anthropic), with the 9 October scope review assisted by Codex
(OpenAI), under the author's direction; the author has not yet read this revision, and no physicist has reviewed
it. The
repository holds every configuration, seed, result and pre-registration:
\url{https://github.com/EmilySmithCreate/RecreationalPhysics}.
\end{acknowledgments}

\begin{thebibliography}{9}
\bibitem{repo} \url{https://github.com/EmilySmithCreate/RecreationalPhysics}: \path{docs/papers/series_plan.md};
PREREGISTRATION T21, T26, T27, T34, T42, T47 (part B), T57; ASSUMPTIONS O20, O21, O22, O25, O36, O41, O43, O48, O49,
O50, O51, O53, O55, O58, O64, O66, O69, O72, O75, O80, O82, O96, O103, O106, O114, Q12, Q18, Q20, Q21, Q22; VISION
Updates 5, 19, 23, 24, 25, 27, 36, 47, 50, 51, 54.
\bibitem{paper1} E. Smith, The curled torus burps: activated escape and conversion in a graph model of emergent
geometry, manuscript (2026), \path{docs/papers/01_curled_torus/}.
\bibitem{paper2} E. Smith, What the change leaves behind, draft (2026), \path{docs/papers/02_relic/}.
\bibitem{paper5} E. Smith, How curled directions open, three and four, draft (2026),
\path{docs/papers/05_six_links/}.
\bibitem{T25} C. A. Trugenberger, Networks as the fundamental constituents of the universe, J. Phys. Complex. 6,
042001 (2025); arXiv:2512.17676.
\bibitem{Smolin06} L. Smolin, The status of cosmological natural selection, arXiv:hep-th/0612185 (2006).
\bibitem{T22} C. A. Trugenberger, Emergent time, cosmological constant and boundary dimension at infinity in
combinatorial quantum gravity, JHEP 04, 019 (2022); arXiv:2112.03778.
\bibitem{Creutz83} M. Creutz, Microcanonical Monte Carlo simulation, Phys. Rev. Lett. 50, 1411 (1983).
\bibitem{DQM25} K. Betre and N. Lewis, Dynamical quantum multigraphs, arXiv:2509.08296 (2025).
\bibitem{paper11} E. Smith, A pull between leftovers, draft (2026), \path{docs/papers/11_gravity/}.
\end{thebibliography}

\end{document}
